Was ist der Anwendungsbereich des VAG
und durch wen wird es ausgeübt?

Aufsichtssubjekte des Versicheraufsichtsgesetz (VAG) sind
alle privat und öffentlich-rechtlichen Versicherungsunternehmen
die Privatversicherungen bertreiben und ihren Sitz in Deutschland haben.
Dies umschließt auch Pensionsfonds und Rückversicherer. Es gehört zum Wirtschatfsverwaltungsrecht.

VR mit anderem Sitzland innerhalb der EWR sind dem Recht ihres Landes unterworfen (§61 VAG).

Die Aufsicht wird von der BaFin übernommen (mit Ausnahme von Bundesland spezifischen Versicherungen oder gesonderten
wie den gesetzlichen [Renten,Krankenver.]).
Die Behörde darf aktiv (materiell) in das Geschäft eingreifen. Dies dient zur Wahrung aller Belange der Versicherten.
Den Aspekt der Erfüllung der Verträge nennt man auch Solvenz oder besondere Finanzaufsicht (§ 294 VAG). Der Rest ist die allg. Rechtsaufsicht.

Das VAG wird durch die Rundschreiben der BaFin konkretisiert. Das größte Mittel ist die Anordnung (§ 298 VAG), was ein Verwaltungsakt
zur Vermeidung oder Beseitung von Misständen ist. Es steht ihnen frei jederzeit Niederlassungen in der EWR aufzusuchen und
Auskünfte und Unterlagen anzufordern zur Einsicht in die Geschäftslage.

Klagen werden vor dem Gericht in Frankfurt eingereicht, die Verwaltung findet in Bonn statt.


Was sind aufsichtsrechtliche Bedingungen zur Gründung einer Versicherung?

Sie ergeben sich in erster Linie aus §§8 VAG. Das Unternehmen bedarf eine bestimmte Rechtsform 
(AG/SE, Versicherungsverein auf Gegenseitigkeit VVaG, Körperschaft des öffentlichen-rechtlichen Rechts).
Nach §15 VAG dürfen keine Versicherungsfremde Geschäfte betrieben werden. Es gilt zudem das Prinzip der Spartentrennung 
(z.B. nicht Lebens und Schadenver.).

Es muss ein Geschäftsplan vorgelegt werden, in dem geklärt wird welche Risiken wie gedeckt werden und ein Rückversicherungskonzept.
Es gibt ebenso Mindestkapitalanforderungen.

Zwei Geschäftsführer sind vorzuweisen, die zuverlässig und fachlich geeignet sind ("fit & proper"). Also ausreichende Erfahrung in 
Versicherungssparte und Führung.

Lebens und Krankver. müssen einen verantwortlichen Aktuar stellen (§§ 141, 156 VAG). Substitutive Krankenver. und Pflichtver. müssen die AVB
vorlegen.


Skizieren Sie den Aufbau von Solvency II


Ziel von Solvency II war eine Harmonisierung der aufsichtsrechtlichen Gesichtspunkte für VR auf europäischer Ebene,
insbesondere Solvabilitätsvorschriften zu einem risikoorrientierten System der Finanzaufsicht weiterzuentwickeln. 
Zum Beispiel die Bewertung von Vermögen und Verbindlichkeiten durch Ihre Marktwerte.

Es besteht aus einem 3 Säulen Modell.

1. Säule quantitative Fragestellungen
- Bewertung Aktiva & Passiva
- vert. Rückstellungen
- Eigenmittel
- SCR & MCR
Berechnung via Standard Formel oder internem Model
2. Qualitative Anforderungen 
- Risikostrategie
- Aufbau/Ablauf organisation
- internes Steuerung/Controlling
- interne Revision
Grundsatz doppelte Proportionalität
(Bei Regulierung ist das Risikoprofil des jeweiligen Unternehmens zu berücksichtigen)
- Supervisory Review Process
3. Berichtspflichten
- Bilanz, GuV, Jahresabschluss
- ORSA
- Capital Add Ons
- SCR, MCR
etc. 

Durchgesetzt wird de

gilt 


